% ═══════════════════════════════════════════════════════════════════════════════
\section{Einleitung}
\label{sec:einleitung}
% ═══════════════════════════════════════════════════════════════════════════════

\subsection{Motivation und Problemstellung}
\label{subsec:motivation}

Moderne Softwaresysteme sind in hohem Maße polyglott: Unterschiedliche Dienste
einer verteilten Anwendungslandschaft werden in verschiedenen Programmiersprachen
entwickelt, greifen jedoch häufig auf gemeinsame relationale Datenbanksysteme zu.
Python-basierte Datenpipelines, Java-Microservices und C\#-Backend-Dienste
formulieren dabei konzeptuell identische Datenbankabfragen -- verwenden jedoch
vollständig unterschiedliche Ausdrucksmittel.

Betrachten wir die simple Frage \emph{,,Alle aktiven Nutzer, deren Umsatz im
vergangenen Quartal über 1\,000~Euro lag, sortiert nach Name''}: In Python
(SQLAlchemy), Java (Stream API mit JPA Criteria) und C\#/.NET (LINQ to Entities)
entstehen syntaktisch völlig unterschiedliche Ausdrücke, die jedoch -- bei
korrekter Implementierung -- denselben SQL-Ausdruck erzeugen sollten. In der
Praxis tun sie dies \emph{nicht}: ORMs und Query-Builder unterschiedlicher
Sprachen und Frameworks produzieren inkompatible, unvollständige oder semantisch
abweichende SQL-Ausgaben \cite{ansi_sql_2016,hibernate_docs,ef_core_docs}.

Daraus entstehen konkrete Probleme in der Softwareentwicklung:

\begin{itemize}[leftmargin=1.5cm]
  \item \textbf{Inkonsistente Abfrageergebnisse:} Zwei Dienste, die dieselbe
    Geschäftslogik ausdrücken, erhalten unterschiedliche Daten.
  \item \textbf{Erschwerte Code-Reviews:} Entwickler müssen den mentalen Overhead
    betreiben, Query-Ausdrücke verschiedener Sprachen zu verstehen und zu vergleichen.
  \item \textbf{Migrations-Risiken:} Bei der Portierung von Java- nach Python-Diensten
    oder umgekehrt kann keine mechanische Äquivalenzprüfung erfolgen.
  \item \textbf{Fehlende Testbarkeit:} Ohne definiertes Normalformat ist die
    automatisierte Verifikation der Query-Äquivalenz nicht möglich.
\end{itemize}

\subsection{Zielsetzung des Standards}
\label{subsec:ziel}

Der \textbf{Unified Query Translation Layer (UQTL)} adressiert diese Problematik
durch die Definition eines umfassenden technischen Standards, der folgende
Ziele verfolgt:

\begin{enumerate}[leftmargin=1.5cm]
  \item \textbf{Semantische Äquivalenzgarantie:} Ausdrücke derselben logischen
    Abfrage in Python, Java und C\# werden stets in dasselbe normalisierte
    SQL-Statement übersetzt.
  \item \textbf{Intermediäre Darstellung:} Ein sprachunabhängiger
    \emph{Intermediate Query Tree (IQT)} dient als gemeinsame Zwischenrepräsentation.
  \item \textbf{Vollständige Operatorabdeckung:} Alle wesentlichen
    Query-Operatoren (Filter, Projektion, Aggregation, Sortierung, Joins,
    Fensterfunktionen, Unterabfragen) sind standardisiert.
  \item \textbf{Erweiterbarkeit:} Der Standard definiert Erweiterungspunkte
    für datenbankspezifische Dialekte (PostgreSQL, MySQL, Oracle, MSSQL).
  \item \textbf{Konformitätsstufen:} Vier Konformitätsstufen (CL-1 bis CL-4)
    erlauben schrittweise Implementierungen.
\end{enumerate}

\begin{infobox}
\textbf{Kerndefinition des Standards:}

Für jede Menge logisch äquivalenter Query-Ausdrücke $Q_{py}$, $Q_{java}$,
$Q_{cs}$ aus Python, Java bzw.\ C\# gilt unter UQTL:
\[
  \mathrm{SQL}(Q_{py}) \;\equiv_\sigma\; \mathrm{SQL}(Q_{java})
  \;\equiv_\sigma\; \mathrm{SQL}(Q_{cs})
\]
wobei $\equiv_\sigma$ die semantische SQL-Äquivalenz unter einer gegebenen
Datenbankschema-Instanz $\sigma$ bezeichnet.
\end{infobox}

\subsection{Abgrenzung}
\label{subsec:abgrenzung}

Der UQTL-Standard ist \emph{kein} vollständiges ORM-Framework und definiert
\emph{keine} Laufzeitbibliothek. Er spezifiziert:

\begin{itemize}[leftmargin=1.5cm]
  \item Die formale Syntax und Semantik des Intermediate Query Tree (IQT)
  \item Ãœbersetzungsregeln von den drei Quellsprachen in den IQT
  \item Ãœbersetzungsregeln vom IQT nach ANSI SQL:2016
  \item Normalisierungsregeln für eindeutige SQL-Ausgaben
  \item Konformitätsanforderungen für Implementierungen
\end{itemize}

Explizit nicht Teil des Standards sind: Datenbankzugriff, Verbindungsverwaltung,
Entity-Mapping, Caching-Strategien sowie dialektspezifische Optimierungen
(diese können als optionale Erweiterungen definiert werden).

\subsection{Aufbau der Arbeit}
\label{subsec:aufbau}

\ref{sec:grundlagen} legt die theoretischen Grundlagen dar:
relationale Algebra, Query-Abstraktionen und bestehende Ansätze.
\ref{sec:iqt} definiert den Intermediate Query Tree formal.
\ref{sec:operatoren} spezifiziert alle Ãœbersetzungsregeln.
\ref{sec:beispiele} demonstriert den Standard anhand umfangreicher
Praxisbeispiele.
\ref{sec:konformitaet} definiert die Konformitätsstufen.
\ref{sec:implementierung} beschreibt Referenzimplementierungen.
\ref{sec:diskussion} diskutiert Limitierungen und Erweiterungen.
\ref{sec:zusammenfassung} fasst die Ergebnisse zusammen.

